\documentclass[11pt,letterpaper]{article}
\usepackage[T1]{fontenc}
\usepackage[utf8]{inputenc}
\usepackage{amsmath,amsthm,mathtools}
\usepackage{newtxtext,newtxmath}
\usepackage[margin=1in,headheight=15pt]{geometry}
\usepackage{microtype}
\usepackage{booktabs,longtable,array}
\usepackage[dvipsnames]{xcolor}
\usepackage{enumitem}
\usepackage{fancyhdr}
\usepackage[most]{tcolorbox}
\usepackage{listings}
\usepackage{hyperref}
\definecolor{Ink}{HTML}{17384A}
\definecolor{Accent}{HTML}{147D86}
\definecolor{Pale}{HTML}{EDF5F6}
\hypersetup{colorlinks=true,linkcolor=Ink,citecolor=Accent,urlcolor=Accent,
 pdftitle={Prime-power congruences for compositional tree series},
 pdfauthor={Research report prepared with ChatGPT}}
\pagestyle{fancy}
\fancyhf{}
\fancyhead[L]{\small\sffamily Compositional tree series}
\fancyhead[R]{\small\sffamily Prime-power congruences}
\fancyfoot[C]{\thepage}
\renewcommand{\headrulewidth}{0.3pt}
\setlength{\parskip}{0.35em}
\setlength{\parindent}{1.2em}
\setcounter{tocdepth}{2}
\newtheorem{theorem}{Theorem}[section]
\newtheorem{lemma}[theorem]{Lemma}
\newtheorem{proposition}[theorem]{Proposition}
\newtheorem{corollary}[theorem]{Corollary}
\theoremstyle{definition}
\newtheorem{definition}[theorem]{Definition}
\newtheorem{example}[theorem]{Example}
\theoremstyle{remark}
\newtheorem{remark}[theorem]{Remark}
\newcommand{\N}{\mathbb N}
\newcommand{\Z}{\mathbb Z}
\newcommand{\Q}{\mathbb Q}
\newcommand{\F}{\mathbb F}
\newcommand{\A}{A_\ell}
\newcommand{\al}[1]{a_\ell(#1)}
\newcommand{\itA}[1]{A_\ell^{\circ #1}}
\newcommand{\T}{\mathcal T}
\newcommand{\Fix}{\operatorname{Fix}}
\newcommand{\Aut}{\operatorname{Aut}}
\newcommand{\vp}{\nu_p}
\newcommand{\ord}{\operatorname{ord}_x}
\newcommand{\lcm}{\operatorname{lcm}}
\newcommand{\ind}{\mathbf 1}
\newcommand{\fall}[2]{(#1)_{\underline{#2}}}
\newcommand{\oeis}[1]{\href{https://oeis.org/#1}{\texttt{#1}}}
\newtcolorbox{resultbox}{colback=Pale,colframe=Accent,boxrule=0.6pt,
 arc=2pt,left=9pt,right=9pt,top=7pt,bottom=7pt,breakable}
\lstset{basicstyle=\ttfamily\small,breaklines=true,frame=single,
 rulecolor=\color{Accent},columns=fullflexible,showstringspaces=false}
\allowdisplaybreaks
\begin{document}
\begin{titlepage}
\vspace*{0.55in}
{\sffamily\color{Accent}\large DISCRETE MATHEMATICS \quad / \quad RESEARCH REPORT\par}
\vspace{0.35in}
{\sffamily\bfseries\color{Ink}\Huge Prime-power congruences\par
\vspace{0.12in}for compositional\par
\vspace{0.12in}tree series\par}
\vspace{0.4in}
{\Large A counterexample, exact modular corrections,\par
and a finite-core theorem for an OEIS family\par}
\vspace{0.45in}
\begin{resultbox}
For the unique formal series
\[
 A_\ell(x)=x\exp\!\bigl(A_\ell^{\circ\ell}(x)\bigr),\qquad
 A_\ell(x)=\sum_{n\ge1}a_\ell(n)\frac{x^n}{n!},
\]
the conjecture $a_5(n)\equiv n\pmod5$ is false. Its first failure is
$n=23$, and all failures are precisely
\[
 n=20j+3\quad\text{or}\quad n=20j+4,\qquad j\ge1.
\]
The mechanism is a general prime-power congruence for labelled trees with
integer types. It also proves eleven other conjectural assertions in the
four OEIS entries considered here.
\end{resultbox}
\vfill
{\large Prepared for Vladimir Reshetnikov\par}
{\large September 19, 2026\par}
\vspace{0.15in}
{\small Self-contained mathematical proofs; exact-arithmetic verification;\par
LaTeX source, Python programs, and machine-readable certificates.\par}
\end{titlepage}

\begin{abstract}
We study the integer coefficients of the formal functional equation
$A_\ell=x\exp(A_\ell^{\circ\ell})$, where the superscript denotes
composition, not an ordinary power. Four recent OEIS entries attach a
collection of congruence conjectures to the cases $\ell=3,4,5,6$.
We disprove the proposed modulo-$5$ identity for $\ell=5$, determine its
complete infinite set of exceptions, and prove the other eleven
currently conjectural assertions recorded in those entries. The main
method gives considerably more: a coefficient modulo $p^\alpha$ is a
weighted sum over typed rooted trees on fewer than $p^\alpha$ labels.
Consequently every sequence in the family is eventually periodic modulo
$p^\alpha$, with period dividing $p^\alpha(p-1)$ and a uniform onset bound
$p^\alpha$. The ordinary residue generating functions are rational.
We also prove periodicity in the iteration parameter modulo the square
of the modulus, and classify exactly when $a_\ell(n)\equiv n\pmod p$
holds for every $n$: there is no restriction for $p=2$, the condition is
$3\mid\ell$ for $p=3$, and it is $p^2\mid\ell$ for every prime $p\ge5$.
The first counterexample is checked by two independent exact formal-series
algorithms. A reproducible test suite performs 29,122 exact-arithmetic
assertions. The proofs, rather than the finite tests, establish the
infinite statements.
\end{abstract}

\tableofcontents
\newpage

\section{The selected problem and the outcome}
\subsection{What was conjectured}
The starting point is the following family of formal power series:
\begin{equation}\label{eq:main}
 A_\ell(x)=x\exp\!\bigl(A_\ell^{\circ\ell}(x)\bigr),
 \qquad A_\ell(0)=0,\quad A_\ell'(0)=1,
 \qquad \ell\in\Z_{\ge0}.
\end{equation}
Here $A^{\circ k}$ means $k$-fold composition, with $A^{\circ0}(x)=x$.
This distinction is essential: replacing composition by $A(x)^\ell$
would give an entirely different problem.

Write
\begin{equation}\label{eq:coeffdef}
 A_\ell(x)=\sum_{n\ge1}a_\ell(n)\frac{x^n}{n!}.
\end{equation}
The cases $\ell=3,4,5,6$ are OEIS
\oeis{A396803}, \oeis{A396804}, \oeis{A396805}, and
\oeis{A396806}, respectively \cite{oeis3,oeis4,oeis5,oeis6}.
The entries were introduced in June 2026. In the versions inspected for
this report, they contain parity and small-modulus conjectures.
The selected assertion is the especially simple claim
\begin{equation}\label{eq:false}
 a_5(n)\equiv n\pmod5\qquad(n\ge1).
\end{equation}
It is attractive because its definition is exact, its truth can be tested
without numerical approximation, and any counterexample settles its
universal quantifier.

\begin{theorem}[Complete correction of the selected conjecture]\label{thm:headline}
For every $n\ge1$,
\begin{equation}\label{eq:headline}
 a_5(n)\equiv
 \begin{cases}
 2,&n=20j+3\text{ or }n=20j+4\text{ for some }j\ge1,\\
 n,&\text{otherwise}
 \end{cases}\pmod5.
\end{equation}
Thus the first counterexample to \eqref{eq:false} is $n=23$.
The exceptional set has natural density $1/10$.
\end{theorem}
The proof is given in Section~\ref{sec:mod5}. It uses only typed trees with
at most four vertices, not a long extrapolation from computed terms.

For direct comparison with the OEIS table, the exact value is
\begin{equation}\label{eq:exact23}
 a_5(23)=516114659489378430688740267292635767160672734031357.
\end{equation}
It ends in $7$, so its residue modulo $5$ is $2$, whereas $23\equiv3$.
The published b-file already contains this integer \cite{bfile5}.
The contradiction is therefore already latent in the published data;
we do not claim to have discovered a previously unavailable term.
The substantive extension is the structural proof, the exact infinite
correction, and the general prime-power theory developed below.

\subsection{Scope of the resolution}
Table~\ref{tab:status} records the assertions addressed. Repeated parity
assertions are counted separately because they are separately stated in
the entries. There are eleven proofs and one disproof among the twelve
assertions listed.

\begin{table}[ht]
\centering\small
\begin{tabular}{@{}p{0.14\textwidth}p{0.57\textwidth}p{0.19\textwidth}@{}}
\toprule
Entry & Assertion in the inspected entry & Outcome here\\
\midrule
A396803 & Parity agrees with $n$; residue modulo $3$ agrees with $n$. & Both proved.\\
A396804 & Parity; $a(n)\equiv n\pmod4$; residue $3$ never occurs modulo $5$; corrected six-term pattern modulo $3$. & All four proved.\\
A396805 & Parity; eventual pattern $(0,1,0)$ modulo $3$ from $n=3$; $a(n)\equiv n\pmod5$. & First two proved; last disproved.\\
A396806 & Parity; agreement with $n$ modulo $3$ and modulo $6$. & All three proved.\\
\bottomrule
\end{tabular}
\caption{The conjectural assertions resolved in this report. Source entries:
\cite{oeis3,oeis4,oeis5,oeis6}.}\label{tab:status}
\end{table}

A396804 also preserves an older modulo-$3$ proposal, together with a
correction explicitly noting its failure at $n=5$. That older disproof is
credited in the entry to Danial Asaria and is not counted as a new result
here. We prove the corrected pattern.

\subsection{Status and priority}
The source checks were performed on September 19, 2026. The assertions in
Table~\ref{tab:status} remained labelled as conjectures in the inspected
entries. No prior resolution of these exact assertions was identified in
the checks performed. That is a bounded statement about the sources
inspected, not a proof of global publication priority. The use of
permutation actions, labelled-tree enumeration, and integer-valued
polynomials is standard mathematics; all specialized arguments needed
here are supplied in full. The report has not undergone independent peer
review or proof-assistant verification.

\section{Formal existence and division-free computation}\label{sec:formal}
\subsection{Existence, uniqueness, and useful identities}
\begin{proposition}\label{prop:exist}
For every nonnegative integer $\ell$, equation~\eqref{eq:main} has a unique
solution in $x+x^2\Q[[x]]$.
\end{proposition}
\begin{proof}
For series with zero constant coefficient and linear coefficient $1$,
the coefficient of degree $d$ in any fixed compositional iterate depends
only on coefficients through degree $d$ of the original series.
The coefficient of degree $n$ on the right side of \eqref{eq:main}
therefore depends only on coefficients of $A_\ell$ through degree $n-1$.
The linear coefficient is $1$. Inductively, the equation determines every
subsequent coefficient uniquely, and those recursively determined
coefficients satisfy the equation.

Equivalently, start with $A^{[0]}=x$ and repeatedly apply
$B\mapsto x\exp(B^{\circ\ell})$. Agreement through degree $d$ improves to
agreement through degree $d+1$. Consequently the coefficients stabilize
and give the unique solution. This argument is entirely formal and
requires no positive radius of analytic convergence.
\end{proof}

Set $F_k=A_\ell^{\circ k}$ for $k\ge0$. Substitution into
\eqref{eq:main} gives
\[
 F_k=F_{k-1}\exp(F_{\ell+k-1})\qquad(k\ge1).
\]
Multiplying these identities yields
\begin{equation}\label{eq:typesystem}
 F_k=x\exp\left(\sum_{j=\ell}^{\ell+k-1}F_j\right),\qquad k\ge0,
\end{equation}
where the sum is empty when $k=0$, and $F_0=x$. The iteration identities
also appear in the OEIS entries; the derivation is included to keep the
subsequent interpretation self-contained.

The endpoints of the parameter family are helpful checks:
\begin{equation}\label{eq:endpoints}
 A_0(x)=xe^x,\quad a_0(n)=n;
 \qquad A_1(x)=xe^{A_1(x)}.
\end{equation}
For $\ell=1$ the tree interpretation below is the usual enumeration of
rooted labelled trees. Lagrange inversion, or Cayley's formula, gives
$a_1(n)=n^{n-1}$. That familiar special case is only a check, not an input
to any of the congruence proofs.

\subsection{An exact recurrence}
The standard partial exponential Bell polynomials can be introduced
without divisions. Put $B_{0,0}=1$, $B_{n,0}=0$ for $n>0$, and
\begin{equation}\label{eq:bell}
 B_{n,k}=\sum_{j=1}^{n-k+1}\binom{n-1}{j-1}a_\ell(j)B_{n-j,k-1}
 \qquad(1\le k\le n).
\end{equation}
If $c_{k,n}=n![x^n]F_k(x)$, then
\begin{equation}\label{eq:composition}
 c_{0,n}=\ind_{n=1},\quad c_{1,n}=a_\ell(n),\qquad
 c_{k,n}=\sum_{j=1}^n c_{k-1,j}B_{n,j}\quad(k\ge2).
\end{equation}
The identity follows by expanding a composition into partitions of the
$n$ labels: a block of size $j$ receives one of $a_\ell(j)$ inner
structures. Formula~\eqref{eq:bell} distinguishes the block containing
a particular label.

Let $e_n=n![x^n]\exp(F_\ell(x))$, so $e_0=1$. Differentiating the
exponential, or selecting the block containing a distinguished label,
gives
\begin{equation}\label{eq:expcoeff}
 e_n=\sum_{j=1}^n\binom{n-1}{j-1}c_{\ell,j}e_{n-j},
 \qquad a_\ell(n+1)=(n+1)e_n.
\end{equation}
These formulas form an incremental algorithm. Once $a_\ell(n)$ is known,
compute the $n$th Bell row, then $c_{k,n}$ in increasing $k$, then $e_n$,
and finally $a_\ell(n+1)$.

\begin{corollary}\label{cor:integrality}
All $a_\ell(n)$ and all $c_{k,n}$ are nonnegative integers. Moreover,
$n\mid a_\ell(n)$ for every $n\ge1$.
\end{corollary}
\begin{proof}
All recurrences use only sums, products, and integer binomial
coefficients. The last assertion follows immediately from
$a_\ell(n)=n e_{n-1}$.
\end{proof}

The first coefficients are
\begin{align}\label{eq:firstpolys}
 a_\ell(1)&=1,& a_\ell(2)&=2,\nonumber\\
 a_\ell(3)&=6\ell+3,& a_\ell(4)&=48\ell^2+12\ell+4,\\
 a_\ell(5)&=660\ell^3-90\ell^2+50\ell+5.\nonumber
\end{align}
Although these examples are polynomials with integer coefficients, we
will not assume that this property persists at every degree. The
integer-valued polynomial argument in Section~\ref{sec:parameter} is
specifically designed not to make that unjustified inference.

\section{A labelled-tree interpretation}\label{sec:trees}
\begin{definition}
An \emph{$\ell$-admissible typed rooted tree} is a rooted tree whose
vertices carry distinct labels and nonnegative integer types. A vertex
of type $t$ may have any unordered set of children, but each child's
type $s$ must satisfy
\begin{equation}\label{eq:allowed}
 \ell\le s\le\ell+t-1.
\end{equation}
A vertex of type $0$ is therefore a leaf. Write $\T_{\ell,k}(n)$ for
these trees on a specified set of $n$ labels with root type $k$.
For root type $1$, abbreviate $\T_\ell(n)=\T_{\ell,1}(n)$.
\end{definition}

Labels and types have different roles. Labels identify vertices and are
permuted by group actions. Types are integer decorations preserved by
those actions. Distinct vertices may have equal types.

For fixed $\ell,k,n$, there are finitely many such trees. Along a path,
\eqref{eq:allowed} gives $s\le t+\ell-1$. For $\ell\ge1$, every type is
at most $k+(n-1)(\ell-1)$; for $\ell=0$, types strictly decrease along
edges. Thus all type assignments involved are finite.

\begin{theorem}\label{thm:treeinterpretation}
The exponential generating function of $\T_{\ell,k}$ is $F_k$.
In particular,
\begin{equation}\label{eq:treecount}
 |\T_\ell(n)|=a_\ell(n).
\end{equation}
\end{theorem}
\begin{proof}
Let $G_k$ be the exponential generating function for trees with root type
$k$. Remove the root. Each resulting component is a rooted typed tree
whose root type lies in the finite set
$\{\ell,\ldots,\ell+k-1\}$. The components form an unordered set on
disjoint label subsets. The labelled-set construction gives
\[
 G_k=x\exp\left(\sum_{j=\ell}^{\ell+k-1}G_j\right).
\]
This infinite system has a unique coefficientwise solution: the linear
coefficients are $1$ and, simultaneously for all $k$, a coefficient of
degree $n$ uses only degrees at most $n-1$ on the right. Each particular
calculation uses only finitely many indices. By \eqref{eq:typesystem},
$F_k$ solves the same system. Hence $G_k=F_k$.
\end{proof}

\begin{definition}\label{def:weight}
For a typed tree $T$, define its total type by
\[
 S(T)=\sum_{v\in V(T)} t_v.
\]
For $r\ge1$ and $d\ge0$, define the weighted core count
\begin{equation}\label{eq:Wdef}
 W_{r,d}(\ell)=\sum_{T\in\T_\ell(r)} S(T)^d.
\end{equation}
Throughout, $0^0=1$. This convention is necessary when the number of
additional orbits is zero. In our root-type-$1$ family $S(T)\ge1$, but
residue classes of $S(T)$ can be zero.
\end{definition}

The interpretation already explains $n\mid a_\ell(n)$: every possible
root label gives the same count by relabelling. It also turns
composition, initially a formal-series operation, into the concrete
local rule \eqref{eq:allowed}.

\section{Reduction modulo a prime: only a small core survives}\label{sec:prime}
\begin{theorem}[Prime-core congruence]\label{thm:prime}
Let $p$ be prime and write $n=pq+r$, where $q\ge0$ and $0\le r<p$.
Then
\begin{equation}\label{eq:primecore}
 a_\ell(n)\equiv
 \begin{cases}
 0,&r=0,\\
 W_{r,q}(\ell),&1\le r<p
 \end{cases}\pmod p.
\end{equation}
For $n=0$ set $a_\ell(0)=0$.
\end{theorem}
\begin{proof}
Choose a permutation $\sigma$ of the $n$ labels consisting of $q$
disjoint $p$-cycles and $r$ fixed labels. The cyclic group generated by
$\sigma$ acts on $\T_\ell(n)$. Its orbits have sizes $1$ or $p$, so
\[
 a_\ell(n)\equiv |\Fix(\sigma)|\pmod p.
\]
A fixed tree must have a fixed root label. Thus there are no fixed trees
if $r=0$.

Assume $r>0$. The parent of a fixed vertex must itself be fixed: the
parent is uniquely determined and the tree is preserved by $\sigma$.
Consequently the fixed vertices induce a connected rooted tree $T$ on
the $r$ fixed labels. They inherit an admissible type assignment, so
$T\in\T_\ell(r)$.

Each nonfixed label orbit has size $p$, and its type is constant. All
vertices in such an orbit have parents either in a single nonfixed
orbit or at one fixed vertex. Contract the nonfixed orbits. The result
is a rooted forest attached to $T$. No quotient cycle is possible,
because iterating the parent map would then prevent reaching the root.

An equivariant parent map from one nonfixed $p$-orbit to a different
nonfixed $p$-orbit has exactly $p$ choices. Once one vertex's parent is
chosen, equivariance determines all the others. These choices are
independent for distinct quotient edges. A quotient forest with at
least one such edge therefore contributes a multiple of $p$.

The only remaining case modulo $p$ has every nonfixed orbit directly
attached to a fixed vertex. Its $p$ vertices are leaves. An orbit
attached to a fixed vertex of type $t_v$ has exactly $t_v$ allowable
common leaf types. Summing over the attachment vertex gives $S(T)$
choices for each of the $q$ distinguishable nonfixed orbits. Thus the
contribution of $T$ is $S(T)^q$. Summing over $T$ proves the formula.
For $q=0$, this simply counts the core itself.
\end{proof}

The theorem replaces an arbitrarily large coefficient by a finite sum
on at most $p-1$ labels. It is important that it counts fixed trees,
not arbitrary partitions of the labels: the parent-map structure is
what forces every surviving nonfixed orbit to consist of leaves.

\section{The prime-power theorem}\label{sec:primepower}
A direct repetition of the prime proof is insufficient modulo $p^2$:
a quotient edge contributes a factor $p$, not necessarily $p^2$.
The extra divisibility comes from the number of ways to select the
nonfixed orbits participating in a nontrivial quotient forest.

\begin{lemma}[A forest divisibility estimate]\label{lem:forest}
Fix an admissible core $T$ on $r>0$ fixed labels and consider a
cyclic action of order $p$ on $b>0$ additional,
distinguishable label orbits of size $p$. The total number of invariant
trees with this core and with at least one parent edge between two
nonfixed orbits is divisible by
\[
 p^{\nu_p(b)+1}.
\]
Here ``cyclic action of order $p$'' means that each additional orbit has
size $p$; the statement includes $p=2$.
\end{lemma}
\begin{proof}
As in the prime proof, contract the nonfixed orbits to obtain a rooted
forest attached to the fixed core. Remove all one-vertex components of
this forest. The remaining forest $H$ has $D$ vertices and $e\ge1$
internal edges. Every component of $H$ has at least two vertices. Keep
its types and its attachment vertices in the fixed core as decorations.

For a fixed isomorphism class of decorated $H$, assigning distinct names
from the $b$ available orbits to its vertices can be done in
\[
 \frac{\fall bD}{|\Aut(H)|}
\]
ways. Each internal edge contributes $p$ equivariant parent maps. Each
of the remaining $b-D$ orbits is a one-vertex component and has $S(T)$
attachment-and-type choices. The contribution is therefore
\begin{equation}\label{eq:forestterm}
 \frac{\fall bD}{|\Aut(H)|}\,p^e S(T)^{b-D}.
\end{equation}
When $b<D$ this class contributes zero and may be omitted.

The action of $\Aut(H)$ on its $e$ directed internal edges is faithful.
Indeed, every vertex is an endpoint of such an edge because there are
no one-vertex components. An automorphism fixing every directed edge
fixes every endpoint, hence every vertex. Thus $\Aut(H)$ embeds into the
symmetric group on $e$ objects, and
\[
 \nu_p(|\Aut(H)|)\le\nu_p(e!)\le e-1.
\]
The last inequality follows, for example, from
$\nu_p(e!)=(e-s_p(e))/(p-1)$, where $s_p(e)\ge1$ is the base-$p$ digit
sum; it also follows by counting multiples of $p,p^2,\ldots$ directly.
Since $D\ge1$, the falling factorial $\fall bD$ contains the factor $b$.
The $p$-adic valuation of \eqref{eq:forestterm} is at least
\[
 \nu_p(b)+e-\nu_p(|\Aut(H)|)\ge\nu_p(b)+1.
\]
The quotient in \eqref{eq:forestterm} is an integer: it counts labelings
of the decorated forest. The valuation calculation is valid even
without using this integrality separately. Summing all finitely many
classes proves the result.
\end{proof}



\begin{theorem}[Prime-power core congruence]\label{thm:primepower}
Let $p$ be prime, $\alpha\ge1$, and $m=p^\alpha$. Write
$n=mq+r$ with $q\ge0$ and $0\le r<m$. Then
\begin{equation}\label{eq:powercore}
 a_\ell(n)\equiv
 \begin{cases}
 0,&r=0,\\
 \displaystyle\sum_{T\in\T_\ell(r)}
 S(T)^{p^{\alpha-1}q},&1\le r<m
 \end{cases}\pmod m.
\end{equation}
The same statement holds for the coefficient sequence of $F_k$,
with cores whose root has type $k$.
\end{theorem}
\begin{proof}
The case $q=0$ is just the defining count of the core, so suppose $q>0$.
Choose a permutation $\sigma$ consisting of $q$ cycles of length
$p^\alpha$ and $r$ fixed labels. Let
$\tau=\sigma^{p^{\alpha-1}}$, the generator of its order-$p$ subgroup.
An orbit of the full cyclic group on the set of trees has size $p^j$
for some $0\le j\le\alpha$. If $j<\alpha$, every point of that orbit is
fixed by $\tau$; if $j=\alpha$, no point is fixed by $\tau$.
Therefore
\begin{equation}\label{eq:subgroup}
 |\T_\ell(n)|\equiv|\Fix(\tau)|\pmod{p^\alpha}.
\end{equation}

On the labels, $\tau$ has $r$ fixed points and
\[
 b=p^{\alpha-1}q
\]
nonfixed orbits of size $p$. If $r=0$, a fixed root is impossible.
Otherwise a $\tau$-fixed tree has a core in $\T_\ell(r)$, exactly as
before. By Lemma~\ref{lem:forest}, every contribution having a
nontrivial quotient component is divisible by
$p^{\nu_p(b)+1}$, hence by $p^\alpha$.
The surviving singleton components contribute $S(T)^b$.
Combining this count with \eqref{eq:subgroup} proves the formula.
Nothing in the proof depended on the root type being $1$.
\end{proof}

\subsection{A wider class of tree specifications}
The proof applies more generally when a type $t$ has a finite set of
allowable child types, of cardinality $d(t)$. The root type is fixed,
children form an unrestricted unordered set, and the rule does not
depend on the numerical labels. In this setting replace $S(T)$ by
$\sum_v d(t_v)$. The finite-core formula and its prime-power proof remain
unchanged. Our family is the special case $d(t)=t$ with allowable set
$\{\ell,\ldots,\ell+t-1\}$.

This extension does not include arbitrary degree restrictions, ordered
children, or constraints comparing different labels. Such changes can
invalidate the independent-leaf-orbit count and require a new argument.

\section{Periodicity, rational generating functions, and finite certificates}
\label{sec:periods}
\begin{corollary}[Uniform eventual period]\label{cor:period}
For $m=p^\alpha$,
\begin{equation}\label{eq:period}
 a_\ell(n+m(p-1))\equiv a_\ell(n)\pmod m
 \qquad(n\ge m).
\end{equation}
Thus an eventual period divides $m(p-1)$, uniformly in $\ell$.
\end{corollary}
\begin{proof}
For $n=mq+r$ with $q\ge1$, the exponent in \eqref{eq:powercore} is
$p^{\alpha-1}q$. If $p\nmid S$, Euler's theorem gives
\[
 S^{p^{\alpha-1}(q+p-1)}\equiv
 S^{p^{\alpha-1}q}\pmod{p^\alpha}.
\]
If $p\mid S$, both sides vanish: the exponent is at least
$p^{\alpha-1}\ge\alpha$. The elementary inequality
$2^{\alpha-1}\ge\alpha$ proves the latter bound for all $p\ge2$.
The case $r=0$ is zero on both sides.
\end{proof}

\begin{corollary}[Rational residue generating function]\label{cor:gf}
Let
$R_{\ell,m}(z)=\sum_{n\ge1}(a_\ell(n)\bmod m)z^n$
be regarded as a formal series over $\Z/m\Z$. Then
\begin{equation}\label{eq:coregf}
 R_{\ell,m}(z)=
 \sum_{r=1}^{m-1}z^r
 \sum_{T\in\T_\ell(r)}
 \frac{1}{1-S(T)^{p^{\alpha-1}}z^m}.
\end{equation}
Every denominator has constant coefficient $1$, so the expression is
well defined even though $\Z/m\Z$ need not be a field.
\end{corollary}
\begin{proof}
For each $r$ and $T$, sum \eqref{eq:powercore} over $q\ge0$ as a
geometric series. Only finitely many cores occur.
\end{proof}

If $S(T)$ is divisible by $p$, then $S(T)^{p^{\alpha-1}}=0$ modulo $m$
and its summand is simply $z^r$. Otherwise the same base has order
dividing $p-1$. Thus a common denominator is
$1-z^{m(p-1)}$, with a finite polynomial accounting for the initial
coefficients. In particular, multiplying $R_{\ell,m}$ by that denominator
produces a polynomial of degree at most $m(p-1)+m-1$.

For a general modulus $M=\prod_i p_i^{\alpha_i}$, the Chinese remainder
theorem gives the explicit eventual period bound
\begin{equation}\label{eq:compositeperiod}
 L(M)=\lcm_i\bigl(p_i^{\alpha_i}(p_i-1)\bigr),
 \qquad n\ge\max_i p_i^{\alpha_i}.
\end{equation}
An actual least period may be smaller. The bound is an infinite theorem,
not a pattern guessed from a finite residue table.

\subsection{What constitutes a finite certificate?}
For fixed $\ell$ and $m=p^\alpha$, it is enough to record, for each
$1\le r<m$, the distribution of $S(T)$ modulo $m$ over
$\T_\ell(r)$. Evaluating a coefficient then consists of a finite weighted
power sum. A certificate can be large for a large prime power, but its
size does not depend on the queried index $n$.

Alternatively, after the theorem is established, one may compute the
first $m(p-1)+m-1$ coefficients by the division-free recurrence and reduce
all later indices using \eqref{eq:period}. For fixed modulus this gives
random access to a residue using a finite table and modular reduction of
the index. The large integer $a_\ell(n)$ never needs to be constructed.

\section{Periodicity in the iteration parameter and an exact classification}
\label{sec:parameter}
\subsection{Weighted cores are integer-valued polynomials}
\begin{lemma}\label{lem:polynomial}
For $r\ge2$ and $d\ge0$, $W_{r,d}(\ell)$ is a polynomial in $\ell$ with
rational coefficients, takes integer values at all nonnegative integers,
and has degree at most $d+r-2$. Moreover,
\begin{equation}\label{eq:Wzero}
 W_{r,d}(0)=r.
\end{equation}
For $r=1$, $W_{1,d}(\ell)=1$.
\end{lemma}
\begin{proof}
Fix an underlying rooted labelled tree on $r$ vertices. Every child of
the root has forced type $\ell$, since the root has type $1$.
If the root has $c\ge1$ children, there are $r-1-c\le r-2$ remaining
variable types. The summand $S^d$ has total degree $d$ in those types and
$\ell$.

Sum variable types from the leaves toward the root. A variable of type
$s$ ranges from $\ell$ to $\ell+t-1$, where $t$ is its parent's type.
Summing a polynomial over such an interval produces a rational
polynomial and raises total degree by at most one. This follows by
expressing powers in a binomial-polynomial basis and using
$\sum_{j=0}^{N-1}\binom jk=\binom N{k+1}$; equivalently it is polynomial
discrete antidifferentiation. After at most $r-2$ sums the degree is at
most $d+r-2$. Summing the finitely many underlying trees preserves this
bound. Integrality at nonnegative $\ell$ follows from the original count.

When $\ell=0$, the root's children all have type $0$ and hence are
leaves. Only rooted stars are possible, one for each of the $r$ root
labels. Each has total type $1$, proving \eqref{eq:Wzero}.
\end{proof}

We need integer-valued, not necessarily integer-coefficient, polynomials.
Every polynomial $P$ of degree at most $D$ taking integer values on
nonnegative integers has a Newton expansion
\begin{equation}\label{eq:newton}
 P(X)=\sum_{j=0}^D c_j\binom Xj,\qquad
 c_j=\Delta^jP(0)\in\Z.
\end{equation}
The expansion follows from finite differences and agrees with $P$ at
all nonnegative integers, hence as a polynomial.

\begin{lemma}\label{lem:newtonperiod}
Let $m=p^\alpha$. An integer-valued polynomial of degree $D<2m$ satisfies
\[
 P(L+m^2)\equiv P(L)\pmod m\qquad(L\ge0).
\]
\end{lemma}
\begin{proof}
For $1\le i\le D$, we have $i<2p^\alpha\le p^{\alpha+1}$, so
$\nu_p(i)\le\alpha$. Since
\[
 \binom{m^2}{i}=\frac{m^2}{i}\binom{m^2-1}{i-1},
\]
its valuation is at least $2\alpha-\nu_p(i)\ge\alpha$.
Vandermonde's identity now gives
\[
 \binom{L+m^2}{j}
 =\sum_{i=0}^j\binom{m^2}{i}\binom L{j-i}
 \equiv\binom Lj\pmod m
\]
for every $j\le D$. Apply \eqref{eq:newton}.
\end{proof}

\begin{theorem}[Parameter periodicity]\label{thm:paramperiod}
For every integer $M\ge2$, every $\ell\ge0$, and every $n\ge1$,
\begin{equation}\label{eq:paramperiod}
 a_{\ell+M^2}(n)\equiv a_\ell(n)\pmod M.
\end{equation}
In particular,
\begin{equation}\label{eq:sufficient}
 M^2\mid\ell\quad\Longrightarrow\quad
 a_\ell(n)\equiv n\pmod M\quad\text{for all }n\ge1.
\end{equation}
\end{theorem}
\begin{proof}
First take $M=m=p^\alpha$. Write $n=mq+r$. If $r=0$, both coefficients
are zero modulo $m$ by divisibility by $n$. If $r=1$, the core is a single
vertex of total type $1$, so the residue is independent of $\ell$.

For $2\le r<m$, the core exponent is zero when $q=0$. When $q>0$, reduce
$q$ to its representative $t\in\{1,\ldots,p-1\}$ modulo $p-1$.
As in the proof of Corollary~\ref{cor:period}, units have the same power
and nonunits vanish, so the residue is
$W_{r,d}(\ell)$ with $d=p^{\alpha-1}t\le\varphi(m)$.
In either case Lemma~\ref{lem:polynomial} gives degree at most
\[
 d+r-2\le\varphi(m)+m-3<2m.
\]
Lemma~\ref{lem:newtonperiod} proves the claimed parameter period $m^2$.
For arbitrary $M$, apply the prime-power result to each factor
$p^\alpha\Vert M$. The shift $M^2$ is a multiple of each $p^{2\alpha}$,
so the Chinese remainder theorem proves \eqref{eq:paramperiod}.
Finally $a_0(n)=n$, and repeated parameter shifts give
\eqref{eq:sufficient}.
\end{proof}

\subsection{Which parameters give the identity modulo a prime?}
\begin{theorem}[Complete prime-modulus classification]\label{thm:classification}
Fix a prime $p$ and a nonnegative integer $\ell$. Then
\[
 a_\ell(n)\equiv n\pmod p\quad\text{for every }n\ge1
\]
holds exactly under the following conditions:
\begin{equation}\label{eq:classification}
 \begin{array}{c|c}
 p&\text{necessary and sufficient condition}\\\hline
 2&\text{no restriction on }\ell\\
 3&3\mid\ell\\
 p\ge5&p^2\mid\ell.
 \end{array}
\end{equation}
\end{theorem}
\begin{proof}
The cases $2$ and $3$ follow from the explicit formulas in
Section~\ref{sec:smallmod}. For $p\ge5$, the value
$a_\ell(3)=6\ell+3$ forces $p\mid\ell$; write $\ell=ph$.
On three labels the core is a star or a path. Section~\ref{sec:cores}
gives
\[
 W_{3,q}(\ell)=3(1+2\ell)^q+
 6\sum_{t=\ell}^{2\ell-1}(1+\ell+t)^q.
\]
Modulo $p$, the sum for $\ell=ph$ contains $h$ full sets of residues.
For $q=p-1$, it is $-h$ because the nonzero residues have $(p-1)$st
power $1$ and zero contributes zero. The prime-core theorem, with
$n=p(p-1)+3$, therefore gives
\begin{equation}\label{eq:generalwitness}
 a_{ph}\bigl(p(p-1)+3\bigr)\equiv3-6h\pmod p.
\end{equation}
Since $p\ge5$, agreement with $n\equiv3$ forces $p\mid h$.
Thus $p^2\mid\ell$ is necessary. It is sufficient by
Theorem~\ref{thm:paramperiod}.
\end{proof}

In fact the same calculation gives an infinite family of witnesses:
\begin{equation}\label{eq:infiniteprimewitness}
 a_{ph}\bigl(p(p-1)j+3\bigr)\equiv3-6h\pmod p
 \qquad(p\ge5,\ j\ge1).
\end{equation}
For $p\nmid h$, every one of these indices violates agreement with $n$.
We do not claim that this progression always begins at the first failure
for every prime; only the case $p=5$ receives a complete exception-set
analysis below.

\section{Small cores and explicit formulas}\label{sec:cores}
The following finite sums are the computational heart of the concrete
applications. They are exact integer identities, before any modular
reduction. Let $I_\ell=\{\ell,\ldots,2\ell-1\}$; it is empty for
$\ell=0$.
\begin{align}
 W_{1,d}(\ell)&=1,\label{eq:W1}\\
 W_{2,d}(\ell)&=2(1+\ell)^d,\label{eq:W2}\\
 W_{3,d}(\ell)&=3(1+2\ell)^d+
 6\sum_{t\in I_\ell}(1+\ell+t)^d.\label{eq:W3}
\end{align}
For four vertices,
\begin{align}
 W_{4,d}(\ell)={}&4(1+3\ell)^d
 +24\sum_{t\in I_\ell}(1+2\ell+t)^d\nonumber\\
 &+12\sum_{s,t\in I_\ell}(1+\ell+s+t)^d\label{eq:W4}\\
 &+24\sum_{t\in I_\ell}\sum_{u=\ell}^{\ell+t-1}
 (1+\ell+t+u)^d.\nonumber
\end{align}

Here is a direct derivation, including the labelled multiplicities.
A two-vertex tree is a path with two choices for its root label; its
child has type $\ell$. On three vertices there are three rooted stars,
with types $(1,\ell,\ell)$, and six rooted paths, with types
$(1,\ell,t)$ for $t\in I_\ell$.

On four vertices the four rooted shapes are as follows. A star has four
choices of root and types $(1,\ell,\ell,\ell)$. A root with two children,
one of which has one child, has $4\cdot3\cdot2=24$ labelled realizations
and types $(1,\ell,\ell,t)$. A root with one child that has two leaves
has $4\cdot3=12$ realizations and types $(1,\ell,s,t)$.
A path has $4!=24$ realizations and types $(1,\ell,t,u)$, with
$t\in I_\ell$ and $\ell\le u\le\ell+t-1$.
These are precisely the four terms in \eqref{eq:W4}.

At $d=0$ the formulas give $1,2,6\ell+3,48\ell^2+12\ell+4$, agreeing
with \eqref{eq:firstpolys}. Independent Pr\"ufer-code enumeration in the
supplied program checks both the multiplicities and the complete total-type
distributions; the displayed derivation establishes them without code.

\section{Complete formulas modulo \texorpdfstring{$2,3,4$}{2,3,4}}
\label{sec:smallmod}
\subsection{Parity for the entire family}
\begin{corollary}\label{cor:parity}
For every $\ell\ge0$ and $n\ge1$,
\[
 a_\ell(n)\equiv n\pmod2.
\]
\end{corollary}
\begin{proof}
If $n$ is even, divisibility by $n$ suffices. If $n=2q+1$, the
prime-core formula has a one-vertex core with total type $1$, giving
$a_\ell(n)\equiv1^q=1$.
\end{proof}

\subsection{The full modulo-3 formula}
\begin{corollary}\label{cor:mod3}
For $q\ge0$,
\begin{equation}\label{eq:mod3}
 a_\ell(3q)\equiv0,\qquad
 a_\ell(3q+1)\equiv1,\qquad
 a_\ell(3q+2)\equiv2(1+\ell)^q\pmod3.
\end{equation}
\end{corollary}
\begin{proof}
Use \eqref{eq:W1}--\eqref{eq:W2} in Theorem~\ref{thm:prime} for $p=3$.
\end{proof}

When $\ell\equiv0\pmod3$, the residues agree with $n$. When
$\ell\equiv1\pmod3$, the pattern from $n=1$ is
\[
 1,2,0,1,1,0,\quad\text{repeating}.
\]
When $\ell\equiv2\pmod3$, the values at $n=1,2$ are $1,2$, and from
$n=3$ onward the pattern is $0,1,0$ repeating. These three cases prove
the relevant assertions for all four OEIS entries. They also show the
necessity of $3\mid\ell$ for agreement with $n$ at every index: take
$q=1$ in the third formula.

\subsection{The full modulo-4 formula}
\begin{corollary}\label{cor:mod4}
For every $\ell\ge0$ and $q\ge0$,
\begin{align}
 a_\ell(4q)&\equiv0,&
 a_\ell(4q+1)&\equiv1,\nonumber\\
 a_\ell(4q+2)&\equiv2(1+\ell)^{2q}\pmod4.\label{eq:mod4a}
\end{align}
For the remaining residue class,
\begin{equation}\label{eq:mod4b}
 a_\ell(3)\equiv3+2\ell\pmod4,
 \qquad
 a_\ell(4q+3)\equiv3+2\left\lceil\frac\ell2\right\rceil\pmod4
 \quad(q\ge1).
\end{equation}
In particular, $a_\ell(n)\equiv n\pmod4$ for all $n\ge1$ if and only if
$4\mid\ell$.
\end{corollary}
\begin{proof}
Apply Theorem~\ref{thm:primepower} with $p=2$, $\alpha=2$. The exponent
is $2q$, so \eqref{eq:W1} and \eqref{eq:W2} give
\eqref{eq:mod4a}. For a three-vertex core, the star term in
\eqref{eq:W3} is $3$ because $1+2\ell$ is odd. For $q\ge1$, an even
integer to the power $2q$ is zero modulo $4$, and an odd integer to that
power is one. The quantity $1+\ell+t$ is odd precisely when $t$ has the
same parity as $\ell$. Exactly $\lceil\ell/2\rceil$ integers in
$I_\ell$ do. Since $6\equiv2\pmod4$, formula~\eqref{eq:mod4b}
follows. Its $q=0$ case is the exact value $6\ell+3$.

If agreement with $n$ holds for every $n$, the index $3$ forces $\ell$
to be even, and the index $7$ then forces $\ell/2$ to be even. Conversely,
$4\mid\ell$ makes every displayed formula agree with $n$.
\end{proof}

The case $\ell=4$ proves the modulo-$4$ conjecture for A396804. Combining
parity with the modulo-$3$ result proves agreement modulo $6$ whenever
$3\mid\ell$, including $\ell=6$ as conjectured in A396806.

\section{The modulo-5 conjecture: all exceptions and a larger family}
\label{sec:mod5}
We prove more than Theorem~\ref{thm:headline} by allowing any parameter
$\ell=5h$.

\begin{lemma}\label{lem:powersums}
For $q\ge0$, put
\[
 U_q=\sum_{s=0}^4s^q\pmod5,
 \qquad \delta_q=\ind_{q>0\text{ and }4\mid q}.
\]
Then $U_q=-\delta_q$ modulo $5$.
\end{lemma}
\begin{proof}
For $q=0$, the convention $0^0=1$ gives $U_0=5=0$ modulo $5$.
For $q>0$, the zero term vanishes. The nonzero residues form a cyclic
group of order $4$; their $q$th powers sum to zero unless $4\mid q$,
in which case the sum is $4=-1$.
For this small prime one can also check the four exponent residues
$q=1,2,3,4$ directly and use $s^4=1$ for $s\ne0$.
\end{proof}

\begin{theorem}[Exact modulo-$5$ formula for parameters divisible by $5$]
\label{thm:multiple5}
Let $h\ge0$. For every $n\ge1$,
\begin{equation}\label{eq:multiple5}
 a_{5h}(n)\equiv
 \begin{cases}
 3-h,&n=20j+3,\ j\ge1,\\
 4-2h,&n=20j+4,\ j\ge1,\\
 n,&\text{otherwise}
 \end{cases}\pmod5.
\end{equation}
\end{theorem}
\begin{proof}
Write $n=5q+r$. The cases $r=0,1,2$ follow at once from the prime-core
formula: the residues are respectively $0,1,2$ because $1+5h\equiv1$.
For $r=3$, formula~\eqref{eq:W3} consists of a star contribution $3$
and a sum over $5h$ consecutive type values. They contain $h$ copies of
each residue modulo $5$. Therefore
\begin{equation}\label{eq:r3five}
 a_{5h}(5q+3)\equiv3+6hU_q\equiv3+hU_q=3-h\delta_q\pmod5.
\end{equation}

For $r=4$, use the four terms in \eqref{eq:W4}. The star contributes $4$.
The root-with-two-children term contributes $24hU_q$. The double sum
with multiplicity $12$ contributes $12\cdot5h^2U_q=0$ modulo $5$.
It remains to evaluate the path term.

Write $t=5h+5j+s$, where $0\le j<h$ and $0\le s<5$. The inner interval
for $u$ has $h+j$ complete residue blocks, followed by residues
$0,\ldots,s-1$. Because the total type is congruent to $1+s+u$, the
inner sum equals
\[
 (h+j)U_q+\sum_{v=0}^{s-1}(1+s+v)^q\pmod5.
\]
After summing over $s$, the complete-block part is a multiple of $5$.
Thus the entire double sum is
\[
 hV_q,\qquad
 V_q=\sum_{0\le v<s\le4}(1+s+v)^q.
\]
The summand is symmetric in $s,v$. The sum over all ordered pairs is
$5U_q=0$, and the diagonal sum is $U_q$ because $s\mapsto1+2s$ permutes
the residues. Hence
\[
 2V_q=-U_q,\qquad V_q=-\tfrac12U_q\pmod5.
\]
Combining the four shapes gives
\begin{equation}\label{eq:r4five}
 a_{5h}(5q+4)\equiv4+24hU_q-12hU_q
 =4+12hU_q=4-2h\delta_q\pmod5.
\end{equation}
Equations~\eqref{eq:r3five}--\eqref{eq:r4five} are exactly
\eqref{eq:multiple5}, since $\delta_q=1$ precisely for $q=4j$ with $j\ge1$.
\end{proof}

Taking $h=1$ proves Theorem~\ref{thm:headline}. The first exceptional
indices are
\[
 23,24,43,44,63,64,83,84,103,104,\ldots.
\]
The terms $a_5(1),\ldots,a_5(22)$ really do agree with their indices
modulo $5$; the theorem explains why a short initial check is misleading.

\subsection{A rational generating-function correction}
Let $R_{5h,5}(z)=\sum_{n\ge1}a_{5h}(n)z^n$ over $\F_5$. Since
$\sum_{n\ge1}nz^n=z/(1-z)^2$, the exact residue formula is equivalent to
\begin{equation}\label{eq:correctedgf}
 \boxed{\displaystyle
 R_{5h,5}(z)=\frac{z}{(1-z)^2}
 -h\frac{z^{23}+2z^{24}}{1-z^{20}}
 \quad\text{in }\F_5[[z]].}
\end{equation}
This is an ordinary generating function for the residues, not the
original exponential generating function $A_{5h}$.

For $h=1$, the least eventual period is exactly $20$, and this
$20$-periodicity starts at $n=5$. To see minimality, observe that in the
periodic tail the zero residues occur exactly at multiples of $5$.
Every period is therefore divisible by $5$. Since $20$ is a period, a
least period must be one of $5,10,20$: the least period of a periodic
word divides every period, as follows by Euclidean division of shifts.
Periods $5$ and $10$ both fail, comparing $n=23$ with $n=28$ and $n=33$
and translating these comparisons by arbitrarily large multiples of $20$.
Finally, the comparisons $a_5(3)\not\equiv a_5(23)$ and
$a_5(4)\not\equiv a_5(24)$ prevent an onset before $5$.

There are exactly two exceptional residue classes in every block of
length $20$ after the initial block. Their natural density is therefore
$2/20=1/10$. In a full periodic block, residues $0,1,2,3,4$ occur with
frequencies $4,4,6,3,3$, respectively.

\subsection{An immediate repaired conjecture}
Within the subfamily $\ell=5h$, agreement with $n$ modulo $5$ at every
index holds exactly when $5\mid h$, or equivalently $25\mid\ell$.
The general prime classification proves that no other parameters work.
Thus replacing the proposed parameter $5$ by a multiple of $25$ repairs
the identity; the original modulus and its strikingly simple right side
can be retained.

\section{The missing residue for A396804}\label{sec:avoidance}
The neighboring conjecture that $a_4(n)$ is never $3$ modulo $5$ is true.
We give an explicit certificate for every index.

For $1\le r\le4$, let $N_r(s)$ be the number of cores in $\T_4(r)$ with
$S(T)\equiv s\pmod5$. Equations~\eqref{eq:W1}--\eqref{eq:W4}, or their
shape-by-shape derivation, give the exact counts in
Table~\ref{tab:weights4}.
\begin{table}[ht]
\centering
\begin{tabular}{@{}rrrrrrr@{}}
\toprule
$r$ & $N_r(0)$ & $N_r(1)$ & $N_r(2)$ & $N_r(3)$ & $N_r(4)$ & Total\\
\midrule
1&0&1&0&0&0&1\\
2&2&0&0&0&0&2\\
3&6&6&6&0&9&27\\
4&180&192&132&160&156&820\\
\bottomrule
\end{tabular}
\caption{Exact total-type frequencies for parameter $\ell=4$.}
\label{tab:weights4}
\end{table}
For example, the row $r=3$ has three stars of total type $9$, and six
paths for each $t=4,5,6,7$, with total type $5+t$.
For $r=4$, the following four contribution vectors provide a compact
manual certificate for the less immediate last row:
\begin{equation}\label{eq:shape4vectors}
\begin{array}{c|rrrrr}
\text{shape}&s=0&s=1&s=2&s=3&s=4\\\hline
\text{star}&0&0&0&4&0\\
\text{root with two children}&24&24&0&24&24\\
\text{fork below root}&36&48&36&36&36\\
\text{path}&120&120&96&96&96
\end{array}
\end{equation}
Each entry follows by substituting $t,s\in\{4,5,6,7\}$ in the
corresponding term of \eqref{eq:W4}. Their column sums give
Table~\ref{tab:weights4}.

By the prime-core theorem,
\[
 a_4(5q+r)\equiv\sum_{s=0}^4N_r(s)s^q\pmod5.
\]
For $q>0$, the answer depends only on $q\bmod4$. The resulting complete
table is
\begin{equation}\label{eq:avoidtable}
\begin{array}{c|c|rrrr}
 &q=0&q\equiv1&q\equiv2&q\equiv3&q\equiv0\pmod4,\ q>0\\\hline
r=0&0&0&0&0&0\\
r=1&1&1&1&1&1\\
r=2&2&0&0&0&0\\
r=3&2&4&4&0&1\\
r=4&0&0&1&2&0
\end{array}
\end{equation}
No entry equals $3$.
\begin{corollary}\label{cor:avoidance}
For all $n\ge1$, $a_4(n)\not\equiv3\pmod5$.
\end{corollary}
This proves the residue-avoidance conjecture without assuming any of the
other conjectures in A396804.

For comparison, the exact core frequencies for $\ell=5$ are
\begin{equation}\label{eq:weights5}
\begin{array}{c|rrrrr}
 r&s=0&s=1&s=2&s=3&s=4\\\hline
1&0&1&0&0&0\\
2&0&2&0&0&0\\
3&6&9&6&6&6\\
4&252&256&252&252&252.
\end{array}
\end{equation}
Modulo $5$, these yield $W_{3,q}=3+U_q$ and $W_{4,q}=4+2U_q$,
providing a second small certificate for the corrected modulo-$5$
formula in the case $h=1$.

\section{Verification, algorithms, and reproducibility}\label{sec:verification}
\subsection{What the computation establishes}
The supplied test suite completed \textbf{29,122 exact-arithmetic
assertions}. Its broad modular grid covers $0\le\ell\le12$ and
$1\le n\le260$. The checks include the following distinct mechanisms:

\begin{enumerate}[leftmargin=*,itemsep=0.2em]
\item The coefficient $a_5(23)$ and all earlier coefficients agree between
an incremental Bell-polynomial recurrence and an independent ordinary
power-series implementation using exact rational arithmetic.
\item For $0\le\ell\le8$ and core sizes at most four, the explicit
shape formulas agree with independent labelled-tree enumeration using
Pr\"ufer codes and direct enumeration of allowable types.
\item The parity, modulo-$3$, modulo-$4$, and modulo-$5$ formulas agree
with direct coefficient calculations on their stated finite test grids.
Prime-power core checks include moduli $4,8,9,25$, in addition to primes.
\item Eventual-period checks and parameter-square-period checks are made
against direct calculations of the original functional equation. The
coefficient routine does not first reduce the parameter or the index by
the theorem it is intended to test.
\end{enumerate}

The exact counterexample is not inferred from a rounded approximation or
from its magnitude. Both independent computations produce precisely the
integer in \eqref{eq:exact23}. The original OEIS b-file supplies an
external third comparison for that value.

The finite checks do not prove the infinite statements. For example,
testing periodicity through $260$ cannot exclude a later failure. That
exclusion is supplied by the cyclic-group and forest-divisibility proofs.
Conversely, the independent implementations help detect transcription,
indexing, or multiplicity errors in the presentation.

\subsection{Algorithmic details}
The main coefficient routine implements
\eqref{eq:bell}--\eqref{eq:expcoeff}. For a truncation through $N$, the
operation count is
\[
 O(N^3+\ell N^2),
\]
and the storage is $O(N^2+\ell N)$. These are arithmetic-operation
bounds, not bit-complexity bounds: unrestricted integer coefficients grow
rapidly. For a fixed modulus, reduction after each sum keeps all stored
coefficients bounded by that modulus.

The second routine stores ordinary coefficients as Python
\texttt{Fraction} objects. It uses ordinary polynomial convolution,
Horner composition, and the recurrence
\[
 [x^n]e^{H(x)}=\frac1n\sum_{j=1}^n j[x^j]H(x)\,[x^{n-j}]e^{H(x)}.
\]
It then iterates the original functional equation until the requested
truncation has stabilized. This is slower but avoids sharing the Bell
polynomial implementation with the main computation.

The direct core enumerator decodes every Pr\"ufer word of length $r-2$,
chooses every possible root, orients the resulting labelled tree, and
assigns types recursively from the root. It is deliberately restricted
to $r\le5$ in the public interface, because its purpose is independent
verification of small certificates rather than large-scale enumeration.

\subsection{Pitfalls avoided}
An exponential generating function should not be reduced modulo a prime
by treating $x^n/n!$ as an ordinary field coefficient when $p\mid n!$.
Our modular algorithm instead reduces the integer sequences and uses
only integer polynomial recurrences.

Similarly, using only a prime-order action would not justify the
prime-power statement. The passage from the full cyclic group to its
order-$p$ subgroup, followed by the valuation estimate on decorated
forests, is essential. Finally, the cases $q=0$ must be kept separate
when powers of zero residues occur. Confusing $0^0$ with $0$ would erase
the initial exceptions responsible for the first false conjecture.

\subsection{Archive contents and commands}
The archive is self-contained apart from a standard Python installation
and a LaTeX distribution. No third-party Python packages are required.
\begin{lstlisting}
python3 code/verify.py
python3 code/iterative_series.py 5 40
python3 code/iterative_series.py 5 260 --modulus 5
pdflatex -interaction=nonstopmode -halt-on-error article.tex
pdflatex -interaction=nonstopmode -halt-on-error article.tex
\end{lstlisting}

The \texttt{data/} directory contains exact early coefficients, the
modular grid, core-frequency certificates, general-prime witnesses, and
a JSON verification summary. The README specifies the scope of each
file. Running the verification program regenerates these artifacts.
No checksum files, font files, or external literature PDFs are bundled.

\section{Further questions and conclusion}\label{sec:conclusion}
The selected conjecture has a complete negative answer, and all eleven
other live assertions in Table~\ref{tab:status} have proofs. Several
natural refinements remain worthwhile.

The prime-power period bound is uniform, but it need not be the least
period. Determining the exact eventual period and the shortest initial
segment for each pair $(\ell,p^\alpha)$ is a finite problem after the
core theorem, and a more conceptual classification may be possible.
The parameter-square period is also a bound rather than a universal
minimal period. Modulo $2$ the parameter does not matter at all; modulo
$3$ its period is only $3$.

Theorem~\ref{thm:classification} completely classifies agreement with
$n$ modulo a prime. A comparably sharp classification modulo arbitrary
prime powers would strengthen the sufficient condition
$p^{2\alpha}\mid\ell$. Our proof gives a finite decision procedure for
each fixed modulus but does not claim that this sufficient condition is
necessary.

A separate direction is the real or complex asymptotic growth of the
integer coefficients as $n\to\infty$ for fixed $\ell\ge2$.
No analytic convergence or asymptotic expansion is asserted in this
report. The arithmetic results are formal and combinatorial, and their
validity does not depend on resolving that analytic question.

The central lesson is specific rather than merely methodological.
For this iterative exponential family, congruences of arbitrarily large
coefficients are controlled by a fixed finite collection of small typed
trees. A numerical contradiction at $n=23$ becomes an exact infinite
exception set, rational residue generating functions, and a uniform
prime-power theory.

\appendix
\section{A compact proof of the main counterexample family}
For readers wishing to verify only the selected conjecture's resolution,
the following is a short route through the full proof.
The typed-tree interpretation identifies $a_5(n)$ with rooted labelled
trees having root type $1$ and child types in $\{5,\ldots,t+4\}$ when
the parent has type $t$. Apply a permutation with $q$ disjoint $5$-cycles
and $r$ fixed labels, where $n=5q+r$ and $0\le r<5$.
Modulo $5$, invariant trees with an edge between two nonfixed label
orbits vanish, because such an edge has five possible equivariant
parent maps. Every surviving orbit consists of leaves attached directly
to a fixed-label core. Thus its contribution is the core's total type
to the $q$th power.

The total-type distributions on one through four fixed labels are the
four rows of \eqref{eq:weights5}. Reducing those rows modulo $5$ gives
\[
 \begin{array}{c|rrrrr}
 1&0&1&0&0&0\\
 2&0&2&0&0&0\\
 3&1&4&1&1&1\\
 4&2&1&2&2&2.
 \end{array}
\]
Consequently the residue at $n=5q+r$ is $0,1,2,3+U_q,4+2U_q$ for
$r=0,1,2,3,4$, respectively. Since $U_q=0$ except when $q$ is a positive
multiple of $4$, and $U_q=-1$ in that case, the only failures are
$n=20j+3$ and $20j+4$ with $j\ge1$. At both kinds of index the residue
is $2$. The first index is $23$.

Every count in this argument is supported by an explicit small-tree
classification. Its infinite quantifier comes from the group action,
not from the size of a computational search.

\section{Selected exact data and modular witnesses}
The first twelve terms for $\ell=5$ are
\begin{center}\small
\begin{tabular}{@{}rrr@{}}
\toprule
$n$&$a_5(n)$&$a_5(n)\bmod5$\\\midrule
1&1&1\\
2&2&2\\
3&33&3\\
4&1264&4\\
5&80505&0\\
6&7365456&1\\
7&893006317&2\\
8&136393773888&3\\
9&25351203412449&4\\
10&5592284787086560&0\\
11&1436624546038113861&1\\
12&423481344402695225472&2\\
\bottomrule
\end{tabular}
\end{center}
These initial values agree with the source entry and its b-file
\cite{oeis5,bfile5}; they are independently regenerated in the archive.
The file \texttt{a396805\_first40.csv} extends the exact list through
$n=40$.

The general witness formula \eqref{eq:generalwitness} gives:
\begin{center}
\begin{tabular}{@{}rrrrr@{}}
\toprule
$p$&$\ell$&$n=p(p-1)+3$&$a_\ell(n)\bmod p$&$n\bmod p$\\\midrule
5&5&23&2&3\\
7&7&45&4&3\\
11&11&113&8&3\\
13&13&159&10&3\\
\bottomrule
\end{tabular}
\end{center}
The table is finite corroboration of an already proved formula. Only
for $p=5$ does this report assert that the displayed index is the first
failure and list every exceptional index.

\clearpage
\begin{thebibliography}{9}
\bibitem{oeis3}
Paul D. Hanna and contributors,
\emph{OEIS A396803: E.g.f. satisfies $A(x)=x\exp(A^{\circ3}(x))$},
created June 8, 2026; inspected September 19, 2026.
\url{https://oeis.org/A396803}.

\bibitem{oeis4}
Paul D. Hanna, Danial Asaria, Seiichi Manyama, and contributors,
\emph{OEIS A396804: E.g.f. satisfies $A(x)=x\exp(A^{\circ4}(x))$},
created June 9, 2026; inspected September 19, 2026.
\url{https://oeis.org/A396804}.
The entry records the correction to the earlier modulo-$3$ proposal.

\bibitem{oeis5}
Paul D. Hanna, Seiichi Manyama, and contributors,
\emph{OEIS A396805: E.g.f. satisfies $A(x)=x\exp(A^{\circ5}(x))$},
June 2026; inspected September 19, 2026.
\url{https://oeis.org/A396805}.

\bibitem{oeis6}
Paul D. Hanna, Seiichi Manyama, and contributors,
\emph{OEIS A396806: E.g.f. satisfies $A(x)=x\exp(A^{\circ6}(x))$},
June 2026; inspected September 19, 2026.
\url{https://oeis.org/A396806}.

\bibitem{bfile5}
Paul D. Hanna,
\emph{Table of $n,a(n)$ for OEIS A396805, $1\le n\le200$},
inspected September 19, 2026.
\url{https://oeis.org/A396805/b396805.txt}.
In particular, the entry at index $23$ already contradicts the printed
modulo-$5$ conjecture.
\end{thebibliography}
\end{document}
